\subsection{TBM jamming and monitoring-based response modelling}
TBM jamming is a recurrent problem in difficult or squeezing ground because the surrounding rock may converge toward the machine, increase shield loading and generate frictional resistance along the shield. The literature usually distinguishes machine-area problems such as cutterhead sticking and shield jamming, and more recent data-driven studies also use the broader expression TBM jamming accidents. In the present paper, the case is described as TBM shield jamming because the abnormal response is dominated by shield-side pressure, shield displacement, high thrust demand and loss of advance rather than by a purely cutterhead-local obstruction.

Traditional TBM performance studies estimate penetration, advance rate, thrust, torque or related operational responses from rock mass properties and machine parameters \citep{yagiz2008utilizing,hassanpour2010tbm,armaghani2018prediction}. With the accumulation of monitoring data, sequence models also learn lagged dependence from recent operational histories \citep{hochreiter1997long,cho2014learning}. These studies provide important response-prediction baselines, but their spatial unit is usually a tunnel section, a sample or a feature vector. A monitoring curve can describe how thrust, torque, penetration, advance rate or shield pressure changes with time or chainage, but it does not identify which geological voxels are plausible relation partners for the cutterhead, front shield, middle shield or tail shield. The present study therefore does not replace monitoring-based predictors. It asks whether a spatial relation representation provides component-indexed information about the rock--machine configuration that monitoring curves alone do not encode.

\subsection{Geological prediction and the component-relation gap}
Advance geological prediction methods, including TSP and related geophysical or drilling-based approaches, are used to identify faults, fractured zones, water-rich zones, weak surrounding rock and other geological anomalies before or during tunnel excavation \citep{li2019tunnel}. For TBM interpretation, such data are valuable because they describe spatial heterogeneity that is not contained in the monitoring record itself. Voxelisation is useful because it preserves spatial support, local attribute values and chainage-referenced geological context.

Geological spatial representation alone is nevertheless incomplete for component-level interpretation. A voxel field can show where anomalous ground occurs, but it does not specify whether that anomaly is spatially related to the cutterhead or shield, whether the relation is surface-normal compatible, or whether the voxel remains active at the current excavation step. TSP voxels, TBM surface samples and monitoring records have different spatial or observational supports and should not be collapsed into a single tabular unit without retaining relation information \citep{goodchild1992geographical,longley2005geographic,openshaw1984modifiable}. The missing representation is a relation layer that links geological spatial supports to component-labelled machine-side supports during excavation.

\subsection{Spatial relation graphs for moving excavation scenes}
GIScience is concerned not only with locations but also with geographical objects, spatial relations, temporal change and processes \citep{goodchild1992geographical,worboys2004gis,galton2001spatial}. Spatial relation theory provides concepts for neighbourhood, adjacency, containment, intersection, distance and orientation \citep{egenhofer1991point,wu2008contiguity}. TBM excavation fits this perspective because rock units, machine surfaces and monitoring responses are observed in a moving excavation coordinate, and the relevant relations change as the machine advances.

Recent geospatial knowledge-graph and graph data-model studies show how heterogeneous spatial entities can be integrated for retrieval, reasoning, query and provenance \citep{zhu2026knowwheregraph,xu2025spatiotemporal,li2025question,rahimi2021topology}. Scenario-oriented spatiotemporal knowledge-graph models likewise emphasise linking heterogeneous data through meaningful objects and relations \citep{wang2025urbanphysical}. The present paper shares this entity--relation modelling perspective, but it operates at a finer voxel and surface-node support and targets a moving underground engineering object rather than an urban knowledge service or indoor social network. Rock voxels are geological spatial objects, TBM surface nodes are machine-side spatial objects, and rock--TBM candidate edges encode screened spatial relation support.

Graphs are also a natural formalism for representing irregular spatial objects and their relations. Recent geospatial graph studies emphasise that edges should be related to the spatial process of interest and should not be treated as arbitrary connections \citep{zhu2020spatial,desabbata2023graph,liu2025urban,chen2026usgt}. Spatial interaction models further emphasise that relations are shaped by distance decay, compatibility, direction, accessibility and spatial support \citep{wilson1971spatial,fotheringham1989spatial}. For the rock--TBM problem, candidate relations should depend on active-zone membership, distance decay, surface-normal compatibility and excavation state. This is different from a fully connected graph or a purely learned adjacency matrix. Graph neural networks, heterogeneous graph models and spatiotemporal graph models remain relevant as future extensions for learning on structured data \citep{schlichtkrull2018modeling,zhang2019heterogeneous,wu2020comprehensive,zhou2020graph,li2018diffusion,wu2019graph,zhang2020deep,geng2019spatiotemporal,sanchez2020learning,pfaff2021learning}. Digital-twin frameworks in tunnelling likewise motivate structured links among monitoring, prediction and interpretation modules \citep{latif2023digital,apoji2023shaping,hu2025dataknowledge}. The present study adopts the evaluation spirit of recent spatiotemporal graph studies---comparison with static, pooled or ablated alternatives---but applies it to a spatial representation problem. The target is not higher forecasting accuracy, but evidence that the proposed graph provides relation reconstruction, component specificity, robustness to graph-construction settings and edge-level provenance beyond simpler geological summaries.
